% Generated by the latex-docs skill  |  type: paper  |  theme: academic
% Options below can be switched by hand: \frontpagetrue/false, \toctrue/false
\documentclass[11pt,a4paper]{article}

% ---------------- document options ----------------
\newif\ifexamdoc   \examdocfalse
\newif\iffrontpage \frontpagetrue   % separate cover page
\newif\iftoc       \toctrue   % clickable table of contents
\newif\ifsolutions \solutionstrue   % assignment: show solutions

% ---------------- document metadata ----------------
\newcommand{\doctitle}{Programming Assignment 2}
\newcommand{\docsubtitle}{Data Prefetching and Cache Replacement in ChampSim}
\newcommand{\docauthor}{}
\newcommand{\docaffil}{CS683: Advanced Computer Architecture}
\newcommand{\docdate}{\today}
\newcommand{\dochead}{CS683}
\newcommand{\ifnotempty}[2]{\if\relax\detokenize\expandafter{#1}\relax\else#2\fi}

% ======================================================================
%  CORE PACKAGES  (shared by every theme and document type)
% ======================================================================
\usepackage[utf8]{inputenc}
% T1 + Latin Modern when available; otherwise default Computer Modern (keeps fonts scalable)
\IfFileExists{lmodern.sty}{\usepackage[T1]{fontenc}\usepackage{lmodern}}{}

% ---------- page layout: compact, professional margins ----------
\usepackage[a4paper,left=2cm,right=2cm,top=2.4cm,bottom=2.2cm,
            headheight=15pt,headsep=0.6cm,footskip=1cm]{geometry}

% ---------- math ----------
\usepackage{amsmath,amssymb,amsthm,mathtools}

% ---------- colour, graphics, plots ----------
\usepackage[table,dvipsnames]{xcolor}
\usepackage{graphicx}
\usepackage{tikz}
\usetikzlibrary{calc,positioning,arrows.meta}
\usepackage{pgfplots}
\pgfplotsset{compat=1.17}

% ---------- tables ----------
\usepackage{booktabs,array,tabularx,longtable,multirow}

% ---------- figures & captions ----------
\usepackage{caption,subcaption}
\usepackage{float}
\usepackage[section]{placeins}   % floats never drift past the end of their section
\usepackage{wrapfig}
\usepackage{pdflscape}

% ---------- lists, units, chemistry ----------
\usepackage[shortlabels]{enumitem}
\usepackage{siunitx}
\usepackage[version=4]{mhchem}

% ---------- headings, boxes, header/footer, TOC ----------
\usepackage{titlesec}
\usepackage{needspace}   % keep headings with the text that follows
\usepackage[most]{tcolorbox}
\ifexamdoc\else\usepackage{fancyhdr}\fi
\usepackage[titles]{tocloft}
\usepackage{microtype}

% ======================================================================
%  THEME: ACADEMIC  -- burgundy accents, Palatino, small caps, no boxes
% ======================================================================
\IfFileExists{mathpazo.sty}{\usepackage[sc]{mathpazo}}{}
\IfFileExists{tgpagella.sty}{\usepackage{tgpagella}}{}
\linespread{1.04}
\definecolor{accent}{RGB}{122,31,43}       % burgundy: theorems, headings
\definecolor{accentdark}{RGB}{90,20,30}
\definecolor{accentB}{RGB}{52,61,99}       % ink blue: definitions
\definecolor{accentC}{RGB}{86,104,48}      % olive: examples
\definecolor{muted}{RGB}{100,96,92}
\colorlet{tablehead}{accent!10}
\colorlet{tableheadtext}{accentdark}
\colorlet{tablerow}{accent!3}
\colorlet{tablerule}{accentdark}

% ---------- headings ----------
\titleformat{\section}{\Large\scshape\color{accent}}{\thesection.}{0.6em}{}
\titleformat{\subsection}{\large\itshape\color{accentdark}}{\thesubsection}{0.6em}{}
\titleformat{\subsubsection}{\normalsize\itshape\color{accentdark}}{\thesubsubsection}{0.6em}{}
\titleformat{\paragraph}[runin]{\normalsize\scshape\color{accentdark}}{}{0pt}{}
\titlespacing*{\section}{0pt}{18pt plus 3pt minus 2pt}{8pt}
\titlespacing*{\subsection}{0pt}{12pt plus 2pt minus 2pt}{5pt}

% ---------- theorem look: coloured small-caps heads, no boxes ----------
\newcommand{\thmheadfont}{\scshape\bfseries}
\newcommand{\thmskip}{8pt plus 2pt minus 2pt}
\newcommand{\themeboxes}{}

% ---------- header / footer ----------
\ifexamdoc\else
\renewcommand{\sectionmark}[1]{\markright{#1}}
\fancypagestyle{themed}{\fancyhf{}%
  \fancyhead[L]{\small\scshape\color{muted}\dochead}%
  \fancyhead[R]{\small\itshape\color{muted}\nouppercase{\rightmark}}%
  \fancyfoot[C]{\small\color{accent}--\,\thepage\,--}%
  \renewcommand{\headrulewidth}{0.3pt}%
  \renewcommand{\headrule}{{\color{accent!50}\hrule height\headrulewidth width\headwidth}}}
\fancypagestyle{plain}{\fancyhf{}\fancyfoot[C]{\small\color{accent}--\,\thepage\,--}%
  \renewcommand{\headrulewidth}{0pt}}
\pagestyle{themed}
\fi

% ---------- title block (no front page) ----------
\newcommand{\makesimpletitle}{%
  \begin{center}
    {\color{accent}\rule{0.12\linewidth}{0.6pt}}\par\vspace{8pt}
    {\LARGE\scshape\color{accent}\doctitle\par}
    \ifnotempty{\docsubtitle}{\vspace{5pt}{\large\itshape\docsubtitle\par}}
    \vspace{8pt}
    {\normalsize\docauthor\par}
    {\small\itshape\color{muted}\docaffil\ifnotempty{\docdate}{\ifnotempty{\docaffil}{, }\docdate}\par}
    \vspace{6pt}{\color{accent}\rule{0.12\linewidth}{0.6pt}}
  \end{center}
  \vspace{8pt}
  \thispagestyle{plain}}

% ---------- front page ----------
\newcommand{\makefrontpage}{%
  \begin{titlepage}
    \centering
    \vspace*{1.5cm}
    {\large\scshape\color{muted}\docaffil\par}
    \vspace{1.2cm}
    {\color{accent}\rule{\linewidth}{1.2pt}}\par\vspace{-8pt}
    {\color{accent}\rule{\linewidth}{0.4pt}}\par
    \vspace{1.1cm}
    {\fontsize{26}{32}\selectfont\scshape\color{accent}\doctitle\par}
    \ifnotempty{\docsubtitle}{\vspace{12pt}{\Large\itshape\docsubtitle\par}}
    \vspace{1.1cm}
    {\color{accent}\rule{\linewidth}{0.4pt}}\par\vspace{-8pt}
    {\color{accent}\rule{\linewidth}{1.2pt}}\par
    \vspace{2.2cm}
    \ifnotempty{\docauthor}{{\large\itshape by}\par\vspace{6pt}{\Large\scshape\docauthor\par}}
    \frontpageextra
    \vfill
    {\color{accent}\large$\diamond$}\par\vspace{6pt}
    {\large\docdate\par}
  \end{titlepage}}

\newcommand{\themeextras}{%
  \renewcommand{\cfttoctitlefont}{\Large\scshape\color{accent}}
  \renewcommand{\cftsecfont}{\scshape\color{accentdark}}
  \renewcommand{\cftsecpagefont}{\color{accentdark}}}

% ======================================================================
%  CORE SETUP  (runs after the theme; uses the theme's colours/fonts)
% ======================================================================

% ---------- shared macros ----------
\newcommand{\R}{\mathbb{R}}\newcommand{\N}{\mathbb{N}}\newcommand{\Z}{\mathbb{Z}}
\newcommand{\Q}{\mathbb{Q}}\newcommand{\C}{\mathbb{C}}
\DeclarePairedDelimiter{\abs}{\lvert}{\rvert}
\DeclarePairedDelimiter{\norm}{\lVert}{\rVert}
\newcommand{\dd}{\mathop{}\!\mathrm{d}}
\newcommand{\vect}[1]{\mathbf{#1}}
\newcommand{\Prob}{\mathbb{P}}
\newcommand{\E}{\mathbb{E}}
\newcommand{\ind}[1]{\mathbf{1}\{#1\}}
\newcommand{\unclear}[1]{\textcolor{red}{[#1?]}}

% ---------- theorem-like environments (styled by the theme) ----------
\newtheoremstyle{themethm}{\thmskip}{\thmskip}{\itshape}{}{\thmheadfont\color{accent}}{.}{.5em}{}
\newtheoremstyle{themedef}{\thmskip}{\thmskip}{\normalfont}{}{\thmheadfont\color{accentB}}{.}{.5em}{}
\newtheoremstyle{themeex}{\thmskip}{\thmskip}{\normalfont}{}{\thmheadfont\color{accentC}}{.}{.5em}{}
\newtheoremstyle{themerem}{\thmskip}{\thmskip}{\normalfont}{}{\thmheadfont\color{muted}}{.}{.5em}{}
\ifexamdoc\else
\theoremstyle{themethm}
\newtheorem{theorem}{Theorem}[section]
\newtheorem{lemma}[theorem]{Lemma}
\newtheorem{proposition}[theorem]{Proposition}
\newtheorem{corollary}[theorem]{Corollary}
\theoremstyle{themedef}
\newtheorem{definition}[theorem]{Definition}
\theoremstyle{themeex}
\newtheorem{example}[theorem]{Example}
\theoremstyle{themerem}
\newtheorem{remark}[theorem]{Remark}
\themeboxes            % theme hook: wraps the environments above in boxes (or not)
\renewcommand{\qedsymbol}{\textcolor{accent}{\ensuremath{\blacksquare}}}
\fi

% ---------- tables ----------
% Use \tablesetup at the start of every table/longtable (make_table.py does this).
\newcolumntype{L}{>{\raggedright\arraybackslash}X}
\newcolumntype{C}{>{\centering\arraybackslash}X}
\newcommand{\tablesetup}{%
  \arrayrulecolor{tablerule}%
  \setlength{\aboverulesep}{0pt}\setlength{\belowrulesep}{0pt}%
  \renewcommand{\arraystretch}{1.3}%
  \rowcolors{2}{white}{tablerow}}
\newcommand{\colhead}[1]{\textbf{\color{tableheadtext}#1}}
\setlength{\heavyrulewidth}{0.9pt}\setlength{\lightrulewidth}{0.5pt}

% ---------- figures & captions ----------
\graphicspath{{figures/}{images/}{./}}
\captionsetup{font=small,labelfont={bf,color=accentdark},labelsep=period,skip=6pt}
\captionsetup[table]{position=top}
\captionsetup[sub]{font=footnotesize,labelfont={bf,color=accentdark}}
% float placement: allow big floats near their text instead of piling them at the end
\renewcommand{\topfraction}{0.85}\renewcommand{\bottomfraction}{0.6}
\renewcommand{\textfraction}{0.12}\renewcommand{\floatpagefraction}{0.75}
\setcounter{topnumber}{3}\setcounter{bottomnumber}{2}\setcounter{totalnumber}{4}
% placeholder for a figure that is not available yet: \figplaceholder[width]{text}
\newcommand{\figplaceholder}[2][0.6\linewidth]{%
  \fcolorbox{accent!50}{accent!4}{\parbox[c][4cm][c]{#1}{\centering\small\color{accentdark}#2}}}

% ---------- table of contents ----------
\renewcommand{\cftsecfont}{\bfseries\color{accentdark}}
\renewcommand{\cftsecpagefont}{\bfseries\color{accentdark}}
\renewcommand{\cftsecleader}{\cftdotfill{\cftdotsep}}
\renewcommand{\cftsubsecfont}{\color{black!85}}
\setlength{\cftbeforesecskip}{4pt}
\newcommand{\maketoc}{%
  \begingroup\hypersetup{linkcolor=.}\tableofcontents\endgroup
  \iffrontpage\clearpage\else\par\vspace{4pt}{\color{accent!40}\hrule}\vspace{10pt}\fi}

% ---------- links ----------
\usepackage{hyperref}
\hypersetup{colorlinks=true,linkcolor=accentdark,citecolor=accentdark,
  urlcolor=accentB,pdftitle={\doctitle},pdfauthor={\docauthor},bookmarksnumbered=true}
\usepackage[capitalise,noabbrev]{cleveref}
\newcommand{\frontpageextra}{}   % types may add content to the front page (e.g. exam cover)
\themeextras                     % theme hook: last-word overrides (TOC look, captions, ...)

% ======================================================================
%  TYPE: PAPER  -- research paper / report (natbib + BibTeX, references.bib)
% ======================================================================
\usepackage[numbers,sort&compress]{natbib}
\renewenvironment{abstract}{%
  \begin{tcolorbox}[enhanced,frame hidden,colback=accent!4,sharp corners,
     left=14pt,right=14pt,top=8pt,bottom=8pt,after skip=14pt]
  \small\noindent{\thmheadfont\color{accent}Abstract.}\ }{\end{tcolorbox}}

% ---------- custom macros ----------
% ---- packages used by this report ----
\usepackage{listings}

% ---- monochrome override of the academic theme (no coloured accents) ----
\definecolor{accent}{gray}{0.0}
\definecolor{accentdark}{gray}{0.0}
\definecolor{accentB}{gray}{0.15}
\definecolor{accentC}{gray}{0.15}
\definecolor{muted}{gray}{0.40}
\definecolor{codebg}{gray}{0.965}
\colorlet{tablehead}{black!7}
\colorlet{tableheadtext}{black}
\colorlet{tablerow}{black!3}
\colorlet{tablerule}{black}
\hypersetup{linkcolor=black,citecolor=black,urlcolor=black!70,
  pdftitle={CS683 Programming Assignment 2 -- Summary}}

% ---- team details (edit here; used on the front page) ----
\newcommand{\teamname}{TeamName}
\newcommand{\teammembers}{%
  Advay Bapat \quad 23B1225\\
  Aniket Gupta \quad 23B1245\\
  Apoorv Goyal \quad 23B1266\\
  Kundan Maiti \quad 23B1242}

% ---- front page in the style of the PA1 report ----
\renewcommand{\makefrontpage}{%
  \begin{titlepage}
    \centering
    \vspace*{1.5cm}
    {\rule{\linewidth}{1.2pt}}\par\vspace{-8pt}
    {\rule{\linewidth}{0.4pt}}\par
    \vspace{1cm}
    {\fontsize{26}{32}\selectfont\scshape\doctitle\par}
    \vspace{10pt}
    {\Large\itshape\docsubtitle\par}
    \vspace{1cm}
    {\rule{\linewidth}{0.4pt}}\par\vspace{-8pt}
    {\rule{\linewidth}{1.2pt}}\par
    \vspace{2.5cm}
    {\large\scshape Team \teamname\par}
    \vspace{0.5cm}
    {\large\teammembers\par}
    \vspace{0.8cm}
    {\large Department of Electrical Engineering\par}
    {\large Indian Institute of Technology Bombay\par}
    \vfill
    {\large\textbf{\docaffil}\par}
    \vspace{0.3cm}
    {\large\textbf{Date:} \docdate\par}
    \vspace{1cm}
  \end{titlepage}}

% ---- header/footer as in the PA1 report: section left, course right ----
\renewcommand{\sectionmark}[1]{\markright{\thesection\quad #1}}
\fancypagestyle{themed}{\fancyhf{}%
  \fancyhead[L]{\small\color{muted}\nouppercase{\rightmark}}%
  \fancyhead[R]{\small\color{muted}CS683}%
  \fancyfoot[C]{\small\textbf{\thepage}}%
  \renewcommand{\headrulewidth}{0.4pt}%
  \renewcommand{\footrulewidth}{0.4pt}%
  \renewcommand{\headrule}{\hrule height\headrulewidth width\headwidth}}
\fancypagestyle{plain}{\fancyhf{}\fancyfoot[C]{\small\textbf{\thepage}}%
  \renewcommand{\headrulewidth}{0pt}\renewcommand{\footrulewidth}{0pt}}
\pagestyle{themed}

% ---- C++ code listing (monochrome version of the PA1 cppcode box) ----
% usage: \begin{cppcode}[Title (\texttt{file.l1d\_pref})] ... \end{cppcode}
\newtcblisting{cppcode}[1][]{%
  listing only, breakable, enhanced,
  listing options={language=C++,basicstyle=\ttfamily\small,
    keywordstyle=\bfseries,commentstyle=\color{black!55}\itshape,
    stringstyle=\color{black!75},numberstyle=\tiny\color{black!50},
    numbers=left,numbersep=10pt,showstringspaces=false,breaklines=true,
    tabsize=4,columns=fullflexible,keepspaces=true},
  colback=codebg,colframe=black!60,boxrule=0.5pt,arc=2pt,
  left=6pt,right=6pt,top=2pt,bottom=2pt,
  title={#1},fonttitle=\bfseries\footnotesize,coltitle=black,
  attach boxed title to top left={xshift=10pt,yshift=-2.5mm},
  boxed title style={colback=white,colframe=black!60,boxrule=0.5pt,arc=1pt,size=small}}

% ---- key-idea box (monochrome version of the PA1 aestheticquote) ----
\newtcolorbox{keyidea}[1][]{enhanced,breakable,frame hidden,boxrule=0pt,
  colback=codebg,borderline west={2.5pt}{0pt}{black},sharp corners,
  left=10pt,right=10pt,top=6pt,bottom=6pt,#1}

% ---- placeholders ----
% \tbd           : empty cell to be filled with a number
% \todo{text}    : grey reminder of what to write (delete when done)
\newcommand{\tbd}{\textcolor{black!40}{\textnormal{--}}}
\newcommand{\todo}[1]{\textcolor{black!55}{\textit{[#1]}}}

% \resultfig[width]{file-stem}{caption}{label}
% Includes figures/<file-stem>.pdf or .png if it exists, otherwise a placeholder box.
\newcommand{\resultfig}[4][0.75\linewidth]{%
  \begin{figure}[htbp]\centering
  \IfFileExists{figures/#2.pdf}{\includegraphics[width=#1]{#2}}{%
  \IfFileExists{figures/#2.png}{\includegraphics[width=#1]{#2}}{%
  \figplaceholder[#1]{Plot: \texttt{figures/\detokenize{#2}}.pdf / .png}}}%
  \caption{#3}\label{#4}
  \end{figure}}

% ---- notation ----
\newcommand{\geomean}{\textit{Geomean}}

\begin{document}

\iffrontpage\makefrontpage\else\makesimpletitle\fi
\iftoc\maketoc\fi


% ======================================================================
%  HOW TO USE THIS TEMPLATE
%  * \tbd marks a number to fill in; \todo{...} marks text to write.
%  * Put plots in figures/ with the file names shown in the grey boxes
%    (.pdf or .png); they appear automatically.
%  * Speedup = IPC(config) / IPC(baseline). Baseline = no L1D prefetcher,
%    SRRIP/LRU at L2 as stated in Sec. 1. Report the geometric mean.
% ======================================================================

% ==========================================
% OVERVIEW
% ==========================================
\section{Overview and Methodology}\label{sec:overview}

\todo{Two or three sentences: what was implemented in each task and the headline result (best geomean speedup and which configuration achieved it).}

\subsection{Simulation Setup}\label{sec:setup}
All experiments use the ChampSim simulator from the PA2 repository, with
\num{25}M warm-up instructions followed by \num{25}M simulated instructions
on each of the three provided traces (\texttt{Trace1}, \texttt{Trace2},
\texttt{Trace3}). \Cref{tab:config} lists the simulated system.

\begin{table}[htbp]
\centering
\caption{Simulated single-core system.}\label{tab:config}
\tablesetup
\begin{tabular}{ll}
\toprule
\rowcolor{tablehead}\colhead{Component} & \colhead{Configuration} \\
\midrule
Branch predictor & Hashed perceptron \\
L1I cache & 32\,KB, 64 sets $\times$ 8 ways, 4-cycle latency \\
L1D cache & 48\,KB, 64 sets $\times$ 12 ways, 5-cycle latency, 16 MSHRs \\
L2 cache  & 512\,KB, 1024 sets $\times$ 8 ways, 10-cycle latency, 32 MSHRs \\
LLC       & 2\,MB, 2048 sets $\times$ 16 ways, 20-cycle latency \\
Block / page size & 64\,B / 4\,KB \\
Warm-up / simulation & 25M / 25M instructions \\
\bottomrule
\end{tabular}
\end{table}

\subsection{Configurations Evaluated}\label{sec:configs}
\begin{table}[htbp]
\centering
\caption{Binaries built for this report (\texttt{./build\_champsim.sh <L1D prefetcher> <L2C replacement> <name>}).}\label{tab:binaries}
\tablesetup
\begin{tabular}{llll}
\toprule
\rowcolor{tablehead}\colhead{Name} & \colhead{L1D prefetcher} & \colhead{L2C replacement} & \colhead{Used in} \\
\midrule
Baseline        & \texttt{no}        & \texttt{lru}       & All tasks \\
IP-stride       & \texttt{ip\_stride} & \texttt{lru}      & Task 1 \\
SPP             & \texttt{spp}       & \texttt{lru}       & Tasks 1, 3 \\
RRIP            & \texttt{no}        & \texttt{rrip}      & Task 2 \\
I-RRIP          & \texttt{no}        & \texttt{irrip}     & Task 2A \\
Dyn.\ I-RRIP    & \texttt{no}        & \texttt{dyn\_irrip} & Tasks 2B, 3 \\
SPP + Dyn.\ I-RRIP & \texttt{spp}    & \texttt{dyn\_irrip} & Task 3 \\
\bottomrule
\end{tabular}
\end{table}

\subsection{Metrics}\label{sec:metrics}
\begin{align}
\text{Speedup} &= \frac{\text{IPC with optimisation}}{\text{IPC of baseline (no prefetching)}} \label{eq:speedup}\\[4pt]
\text{L1D Coverage} &= \frac{\text{L1D load misses}_{\text{baseline}} - \text{L1D load misses}_{\text{prefetcher}}}{\text{L1D load misses}_{\text{baseline}}} \label{eq:coverage}\\[4pt]
\text{Prefetch Accuracy} &= \frac{\text{useful prefetches}}{\text{useful} + \text{useless prefetches}} \label{eq:accuracy}\\[4pt]
\text{MPKI} &= \frac{\text{load misses}}{\text{instructions}/1000} \label{eq:mpki}
\end{align}
MPKI and prefetch accuracy are taken directly from the ChampSim statistics.
Averages across traces are reported as geometric means for speedup and as
arithmetic means for MPKI, accuracy and coverage.

% ==========================================
% TASK 1
% ==========================================
\section{Task 1: Data Prefetching}\label{sec:task1}

\subsection{Task 1A: Understanding the IP-Stride Prefetcher}\label{sec:task1a}
\todo{Explain \texttt{prefetcher/ip\_stride.l1d\_pref}: the tracker table and its fields, how an entry is indexed by IP, how the stride is learned and confirmed, when and how many prefetches are issued, and how the page boundary is handled.}

\paragraph{Tracker table.} \todo{Number of entries, indexing, replacement of tracker entries.}

\paragraph{Training.} \todo{How the last address and last stride are updated on every access.}

\paragraph{Prefetch issue.} \todo{Condition for issuing prefetches, the degree, and the page-boundary check.}

\paragraph{Limitations.} \todo{Why tracking one last stride per IP cannot capture a repeating sequence of strides}, e.g.\ $(+1,+2,+2,+2,+3)^n$. \todo{This motivates SPP.}

\subsection{Task 1B: Stride Pattern Prefetcher (SPP)}\label{sec:task1b}

\subsubsection{Design Overview}\label{sec:spp-design}
\todo{High-level description of the prefetcher. A block diagram (figures/spp\_diagram) is recommended.}

\resultfig[0.8\linewidth]{spp_diagram}{Organisation of the Stride Pattern Prefetcher at the L1D.}{fig:spp-diagram}

\paragraph{IP table.} \todo{Per-IP state tracked: tag, last address, last stride, confidence, signature, mode.}

\paragraph{Uniform-stride detection.} \todo{How a constant stride is detected and confirmed.}

\paragraph{Fallback to non-uniform mode.} \todo{The confidence counter / state machine that detects that the uniform pattern has broken and switches the IP to signature-based prediction, and how it switches back.}

\paragraph{Signature and pattern table.} \todo{How the signature is built from recent strides (e.g.\ shift-and-XOR of the last $k$ strides), its width, how it indexes the pattern table and what the table predicts.}

\paragraph{Multi-degree prefetching.} \todo{How prefetches beyond the first are generated (stride repeated for uniform mode; signature walked forward for non-uniform mode).}

\paragraph{Page boundary.} Every candidate address is checked against the
4\,KB page of the triggering access, and prefetching stops at the first
candidate that falls outside it. \todo{Confirm / adapt.}

\begin{cppcode}[Key routine (\texttt{prefetcher/spp.l1d\_pref})]
// TODO: paste the core of l1d_prefetcher_operate() here
\end{cppcode}

\subsubsection{Accuracy-Based Throttling}\label{sec:spp-throttle}
\todo{How useful and issued prefetches are counted, the length of an evaluation epoch, and how the degree is changed at the end of each epoch.}

\begin{table}[htbp]
\centering
\caption{Prefetch degree chosen for each accuracy range.}\label{tab:spp-degree}
\tablesetup
\begin{tabular}{lcl}
\toprule
\rowcolor{tablehead}\colhead{Accuracy range} & \colhead{Degree} & \colhead{Reason} \\
\midrule
$\text{acc} \ge \tbd$          & \tbd & \todo{high accuracy: be aggressive} \\
$\tbd \le \text{acc} < \tbd$  & \tbd & \todo{\dots} \\
$\tbd \le \text{acc} < \tbd$  & \tbd & \todo{\dots} \\
$\text{acc} < \tbd$           & \tbd & \todo{low accuracy: avoid pollution} \\
\bottomrule
\end{tabular}
\end{table}

\paragraph{How the thresholds were chosen.}
\todo{Deliverable: what thresholds were used and how they were decided (sweep, per-trace accuracy distribution, etc.). Refer to any sweep data.}

\paragraph{Why these degrees.}
\todo{Deliverable: justify the degree in each accuracy range (\cref{tab:spp-degree}).}

\subsubsection{Storage Overhead}\label{sec:spp-storage}
\todo{Deliverable: exact bit-level breakdown of every table entry. Edit the fields to match your implementation.}

\begin{table}[htbp]
\centering
\caption{Bit-level storage breakdown of the SPP.}\label{tab:spp-storage}
\tablesetup
\begin{tabular}{llrrr}
\toprule
\rowcolor{tablehead}\colhead{Structure} & \colhead{Field} & \colhead{Bits/entry} & \colhead{Entries} & \colhead{Total bits} \\
\midrule
IP table       & Valid                    & 1    & \tbd & \tbd \\
               & IP tag                   & \tbd & \tbd & \tbd \\
               & Last line offset in page & 6    & \tbd & \tbd \\
               & Last stride (signed)     & 7    & \tbd & \tbd \\
               & Stride confidence        & \tbd & \tbd & \tbd \\
               & Mode (uniform / non-uniform) & 1 & \tbd & \tbd \\
               & Signature                & \tbd & \tbd & \tbd \\
               & Replacement (LRU) bits   & \tbd & \tbd & \tbd \\
\midrule
Pattern table  & Predicted stride         & 7    & \tbd & \tbd \\
               & Confidence               & \tbd & \tbd & \tbd \\
\midrule
Throttling     & Issued counter           & \tbd & 1    & \tbd \\
               & Useful counter           & \tbd & 1    & \tbd \\
               & Current degree           & \tbd & 1    & \tbd \\
\midrule
\multicolumn{4}{l}{\textbf{Total}} & \textbf{\tbd\ bits (\tbd\ bytes)} \\
\bottomrule
\end{tabular}
\end{table}

\subsubsection{Results}\label{sec:task1-results}
\Cref{tab:task1-results} compares the baseline, the provided IP-stride
prefetcher and our SPP; \cref{fig:t1-speedup,fig:t1-mpki,fig:t1-acc,fig:t1-cov}
plot the same data.

\begin{table}[htbp]
\centering
\caption{Task 1 results (L1D prefetching, LRU at L2).}\label{tab:task1-results}
\tablesetup
\begin{tabular}{llrrrrr}
\toprule
\rowcolor{tablehead}\colhead{Trace} & \colhead{Config} & \colhead{IPC} & \colhead{Speedup} & \colhead{L1D load MPKI} & \colhead{Accuracy (\%)} & \colhead{Coverage (\%)} \\
\midrule
                          & Baseline  & \tbd & 1.000 & \tbd & -- & -- \\
                         & IP-stride & \tbd & \tbd  & \tbd & \tbd & \tbd \\
\multirow{-3}{*}{Trace 1} & SPP       & \tbd & \tbd  & \tbd & \tbd & \tbd \\
\midrule
                          & Baseline  & \tbd & 1.000 & \tbd & -- & -- \\
                         & IP-stride & \tbd & \tbd  & \tbd & \tbd & \tbd \\
\multirow{-3}{*}{Trace 2} & SPP       & \tbd & \tbd  & \tbd & \tbd & \tbd \\
\midrule
                          & Baseline  & \tbd & 1.000 & \tbd & -- & -- \\
                         & IP-stride & \tbd & \tbd  & \tbd & \tbd & \tbd \\
\multirow{-3}{*}{Trace 3} & SPP       & \tbd & \tbd  & \tbd & \tbd & \tbd \\
\midrule
                             & IP-stride & -- & \tbd & \tbd & \tbd & \tbd \\
\multirow{-2}{*}{Mean} & SPP       & -- & \tbd & \tbd & \tbd & \tbd \\
\bottomrule
\end{tabular}
\end{table}

\resultfig{task1_speedup}{Speedup over the no-prefetching baseline.}{fig:t1-speedup}
\resultfig{task1_l1d_mpki}{L1D load MPKI.}{fig:t1-mpki}
\resultfig{task1_accuracy}{Prefetch accuracy.}{fig:t1-acc}
\resultfig{task1_coverage}{L1D coverage.}{fig:t1-cov}

\subsubsection{Discussion}\label{sec:task1-discussion}
\todo{Explain the results trace by trace: which traces have uniform vs.\ non-uniform patterns, where SPP gains over IP-stride and why, effect of throttling on accuracy, any trace where prefetching hurts.}

% ==========================================
% TASK 2
% ==========================================
\section{Task 2: Cache Replacement Policy}\label{sec:task2}

\subsection{Baseline: SRRIP at the L2}\label{sec:rrip}
The provided \texttt{replacement/rrip.l2c\_repl} implements SRRIP with a
2-bit re-reference prediction value (RRPV, $\text{maxRRPV}=3$):
\begin{itemize}
  \item \textbf{Insertion:} a newly filled line gets $\text{RRPV}=2$ (long re-reference interval).
  \item \textbf{Promotion:} a hit sets $\text{RRPV}=0$ (near-immediate re-reference).
  \item \textbf{Eviction:} the first way with $\text{RRPV}=3$ is chosen; if none exists, all RRPVs in the set are incremented and the search repeats.
\end{itemize}

\subsection{Task 2A: Instruction-Conscious RRIP (I-RRIP)}\label{sec:task2a}

\subsubsection{Interface Change}\label{sec:irrip-interface}
\todo{Describe how \texttt{update\_replacement\_state} was extended with an \texttt{isInstruction} argument: the declaration in \texttt{inc/cache.h}, every call site in \texttt{src/cache.cc}, and how the value is taken from \texttt{packet.instruction}.}

\begin{cppcode}[Extended signature (\texttt{inc/cache.h})]
void l2c_update_replacement_state(uint32_t cpu, uint32_t set, uint32_t way,
        uint64_t full_addr, uint64_t ip, uint64_t victim_addr,
        uint32_t type, uint8_t hit, uint8_t isInstruction);
\end{cppcode}

\subsubsection{Policy Design}\label{sec:irrip-design}
\Cref{tab:irrip-policy} summarises how I-RRIP differs from SRRIP in each
of the three parts of the policy.

\begin{table}[htbp]
\centering
\caption{Insertion, promotion and eviction in SRRIP and I-RRIP.}\label{tab:irrip-policy}
\tablesetup
\begin{tabularx}{\linewidth}{lLL}
\toprule
\rowcolor{tablehead}\colhead{Part} & \colhead{SRRIP} & \colhead{I-RRIP} \\
\midrule
Insertion & RRPV $=2$ for every line & \todo{instruction lines: \dots; data lines: \dots} \\
Promotion & RRPV $=0$ on any hit     & \todo{\dots} \\
Eviction  & First line with RRPV $=3$; age set otherwise & \todo{\dots} \\
\bottomrule
\end{tabularx}
\end{table}

\paragraph{Insertion.} \todo{Rationale for modifying or not modifying insertion.}

\paragraph{Promotion.} \todo{Rationale for modifying or not modifying promotion.}

\paragraph{Eviction.} \todo{Rationale for modifying or not modifying eviction.}

\begin{keyidea}
\textbf{Why this design is the best of those we tried.}
\todo{Deliverable: defend the choice. List the alternatives evaluated (e.g.\ instruction-only promotion, data-first eviction, different insertion RRPVs) and the evidence against each.}
\end{keyidea}

\subsubsection{Results}\label{sec:task2a-results}

\begin{table}[htbp]
\centering
\caption{Task 2A results (no prefetcher).}\label{tab:task2a-results}
\tablesetup
\begin{tabular}{llrrrrr}
\toprule
\rowcolor{tablehead}\colhead{Trace} & \colhead{L2 policy} & \colhead{IPC} & \colhead{Speedup} & \colhead{L2C load MPKI} & \colhead{L2 instr.\ MPKI} & \colhead{L2 data MPKI} \\
\midrule
                          & LRU    & \tbd & 1.000 & \tbd & \tbd & \tbd \\
                         & SRRIP  & \tbd & \tbd  & \tbd & \tbd & \tbd \\
\multirow{-3}{*}{Trace 1} & I-RRIP & \tbd & \tbd  & \tbd & \tbd & \tbd \\
\midrule
                          & LRU    & \tbd & 1.000 & \tbd & \tbd & \tbd \\
                         & SRRIP  & \tbd & \tbd  & \tbd & \tbd & \tbd \\
\multirow{-3}{*}{Trace 2} & I-RRIP & \tbd & \tbd  & \tbd & \tbd & \tbd \\
\midrule
                          & LRU    & \tbd & 1.000 & \tbd & \tbd & \tbd \\
                         & SRRIP  & \tbd & \tbd  & \tbd & \tbd & \tbd \\
\multirow{-3}{*}{Trace 3} & I-RRIP & \tbd & \tbd  & \tbd & \tbd & \tbd \\
\midrule
                             & SRRIP  & -- & \tbd & \tbd & \tbd & \tbd \\
\multirow{-2}{*}{Mean} & I-RRIP & -- & \tbd & \tbd & \tbd & \tbd \\
\bottomrule
\end{tabular}
\end{table}

\resultfig{task2a_speedup}{Task 2A: speedup of SRRIP and I-RRIP.}{fig:t2a-speedup}
\resultfig{task2a_l2_mpki}{Task 2A: L2C, L2 instruction and L2 data load MPKI.}{fig:t2a-mpki}

\paragraph{Discussion.} \todo{Change in L2 instruction MPKI vs.\ data MPKI per trace, and the resulting IPC change. Note any trace where data misses rise enough to hurt.}

\subsection{Task 2B: Dynamic I-RRIP}\label{sec:task2b}

\subsubsection{Set-Dueling Design}\label{sec:dyn-design}
\todo{Describe the leader sets (how many, how they are chosen), what is counted in them (instruction and total accesses and misses), how often the counters are evaluated / reset, and how follower sets pick a policy.}

The decision metric is
\begin{equation}
  r = \frac{\text{instr.\ accesses} + \text{instr.\ misses}}{\text{total accesses} + \text{total misses}},
  \qquad
  \text{policy} =
  \begin{cases}
    \text{I-RRIP} & r > \theta \\
    \text{SRRIP}  & \text{otherwise}
  \end{cases}
  \label{eq:dyn}
\end{equation}
with threshold $\theta = \tbd$. \todo{Edit if you changed the formula or added parameters.}

\begin{cppcode}[Policy selection (\texttt{replacement/dyn\_irrip.l2c\_repl})]
// TODO: paste the leader-set check and decision logic here
\end{cppcode}

\subsubsection{Rationale for the Leader Sets and Decision Logic}\label{sec:dyn-rationale}
\paragraph{Number of leader sets.} \todo{Deliverable: why this number (sampling accuracy vs.\ the cost of running leader sets on the non-preferred policy), and how they are spread over the 1024 sets.}

\paragraph{Decision logic.} \todo{Deliverable: why this metric, why accesses and misses are combined, epoch length / hysteresis.}

\subsubsection{Sensitivity Study}\label{sec:dyn-sensitivity}
\todo{Deliverable: sensitivity to the decision parameters (at least 5 threshold values).}

\begin{table}[htbp]
\centering
\caption{Sensitivity of Dynamic I-RRIP to the threshold $\theta$ (speedup over LRU).}\label{tab:dyn-sensitivity}
\tablesetup
\begin{tabular}{crrrrr}
\toprule
\rowcolor{tablehead}\colhead{$\theta$} & \colhead{Trace 1} & \colhead{Trace 2} & \colhead{Trace 3} & \colhead{Geomean} & \colhead{Mean L2 instr.\ MPKI} \\
\midrule
\tbd & \tbd & \tbd & \tbd & \tbd & \tbd \\
\tbd & \tbd & \tbd & \tbd & \tbd & \tbd \\
\tbd & \tbd & \tbd & \tbd & \tbd & \tbd \\
\tbd & \tbd & \tbd & \tbd & \tbd & \tbd \\
\tbd & \tbd & \tbd & \tbd & \tbd & \tbd \\
\bottomrule
\end{tabular}
\end{table}

\resultfig{task2b_sensitivity}{Speedup of Dynamic I-RRIP against the threshold $\theta$.}{fig:t2b-sens}

\todo{Discuss the trend and the value chosen. Optional: a second sweep over the number of leader sets.}

\subsubsection{Results}\label{sec:task2b-results}

\begin{table}[htbp]
\centering
\caption{Task 2B results (no prefetcher, $\theta = \tbd$).}\label{tab:task2b-results}
\tablesetup
\begin{tabular}{llrrrrr}
\toprule
\rowcolor{tablehead}\colhead{Trace} & \colhead{L2 policy} & \colhead{IPC} & \colhead{Speedup} & \colhead{L2C load MPKI} & \colhead{L2 instr.\ MPKI} & \colhead{L2 data MPKI} \\
\midrule
                          & SRRIP        & \tbd & \tbd & \tbd & \tbd & \tbd \\
                         & I-RRIP       & \tbd & \tbd & \tbd & \tbd & \tbd \\
\multirow{-3}{*}{Trace 1} & Dyn.\ I-RRIP & \tbd & \tbd & \tbd & \tbd & \tbd \\
\midrule
                          & SRRIP        & \tbd & \tbd & \tbd & \tbd & \tbd \\
                         & I-RRIP       & \tbd & \tbd & \tbd & \tbd & \tbd \\
\multirow{-3}{*}{Trace 2} & Dyn.\ I-RRIP & \tbd & \tbd & \tbd & \tbd & \tbd \\
\midrule
                          & SRRIP        & \tbd & \tbd & \tbd & \tbd & \tbd \\
                         & I-RRIP       & \tbd & \tbd & \tbd & \tbd & \tbd \\
\multirow{-3}{*}{Trace 3} & Dyn.\ I-RRIP & \tbd & \tbd & \tbd & \tbd & \tbd \\
\midrule
Mean & Dyn.\ I-RRIP & -- & \tbd & \tbd & \tbd & \tbd \\
\bottomrule
\end{tabular}
\end{table}

\resultfig{task2b_speedup}{Task 2B: speedup of SRRIP, I-RRIP and Dynamic I-RRIP.}{fig:t2b-speedup}
\resultfig{task2b_l2_mpki}{Task 2B: L2C, L2 instruction and L2 data load MPKI.}{fig:t2b-mpki}

\paragraph{Discussion.} \todo{Fraction of time each trace spends in I-RRIP mode, and where dynamic switching recovers data misses that static I-RRIP lost.}

% ==========================================
% TASK 3
% ==========================================
\section{Task 3: SPP with Prefetch-Aware Dynamic I-RRIP}\label{sec:task3}

\subsection{Design}\label{sec:task3-design}
\todo{How \texttt{dyn\_irrip.l2c\_repl} treats lines whose packet type is \texttt{PREFETCH}: insertion RRPV for prefetched lines, what happens on the first demand hit to a prefetched line, whether prefetch hits promote, and any change to eviction priority or to the set-dueling counters.}

\begin{table}[htbp]
\centering
\caption{Treatment of each line class in the prefetch-aware policy.}\label{tab:task3-policy}
\tablesetup
\begin{tabularx}{\linewidth}{lLLL}
\toprule
\rowcolor{tablehead}\colhead{Line class} & \colhead{Insertion} & \colhead{Promotion} & \colhead{Eviction priority} \\
\midrule
Instruction (demand) & \todo{\dots} & \todo{\dots} & \todo{\dots} \\
Data (demand)        & \todo{\dots} & \todo{\dots} & \todo{\dots} \\
Prefetched           & \todo{\dots} & \todo{\dots} & \todo{\dots} \\
\bottomrule
\end{tabularx}
\end{table}

\subsection{Rationale for Handling Prefetched Lines}\label{sec:task3-rationale}
\todo{Deliverable: why prefetched lines are handled this way (e.g.\ inaccurate prefetches should not displace instruction lines; accurate prefetches should survive until their demand use; prefetch fills must not bias the set-dueling counters).}

\subsection{Results}\label{sec:task3-results}

\begin{table}[htbp]
\centering
\caption{Task 3 results: performance and prefetch metrics.}\label{tab:task3-perf}
\tablesetup
\begin{tabular}{llrrrrr}
\toprule
\rowcolor{tablehead}\colhead{Trace} & \colhead{Config} & \colhead{IPC} & \colhead{Speedup} & \colhead{Accuracy (\%)} & \colhead{Coverage (\%)} & \colhead{L1D MPKI} \\
\midrule
                          & SPP                & \tbd & \tbd & \tbd & \tbd & \tbd \\
                         & Dyn.\ I-RRIP       & \tbd & \tbd & --   & --   & \tbd \\
\multirow{-3}{*}{Trace 1} & SPP + Dyn.\ I-RRIP & \tbd & \tbd & \tbd & \tbd & \tbd \\
\midrule
                          & SPP                & \tbd & \tbd & \tbd & \tbd & \tbd \\
                         & Dyn.\ I-RRIP       & \tbd & \tbd & --   & --   & \tbd \\
\multirow{-3}{*}{Trace 2} & SPP + Dyn.\ I-RRIP & \tbd & \tbd & \tbd & \tbd & \tbd \\
\midrule
                          & SPP                & \tbd & \tbd & \tbd & \tbd & \tbd \\
                         & Dyn.\ I-RRIP       & \tbd & \tbd & --   & --   & \tbd \\
\multirow{-3}{*}{Trace 3} & SPP + Dyn.\ I-RRIP & \tbd & \tbd & \tbd & \tbd & \tbd \\
\midrule
Mean & SPP + Dyn.\ I-RRIP & -- & \tbd & \tbd & \tbd & \tbd \\
\bottomrule
\end{tabular}
\end{table}

\begin{table}[htbp]
\centering
\caption{Task 3 results: L2 load MPKI.}\label{tab:task3-l2}
\tablesetup
\begin{tabular}{llrrr}
\toprule
\rowcolor{tablehead}\colhead{Trace} & \colhead{Config} & \colhead{L2C load MPKI} & \colhead{L2 instr.\ MPKI} & \colhead{L2 data MPKI} \\
\midrule
                          & SPP                & \tbd & \tbd & \tbd \\
                         & Dyn.\ I-RRIP       & \tbd & \tbd & \tbd \\
\multirow{-3}{*}{Trace 1} & SPP + Dyn.\ I-RRIP & \tbd & \tbd & \tbd \\
\midrule
                          & SPP                & \tbd & \tbd & \tbd \\
                         & Dyn.\ I-RRIP       & \tbd & \tbd & \tbd \\
\multirow{-3}{*}{Trace 2} & SPP + Dyn.\ I-RRIP & \tbd & \tbd & \tbd \\
\midrule
                          & SPP                & \tbd & \tbd & \tbd \\
                         & Dyn.\ I-RRIP       & \tbd & \tbd & \tbd \\
\multirow{-3}{*}{Trace 3} & SPP + Dyn.\ I-RRIP & \tbd & \tbd & \tbd \\
\bottomrule
\end{tabular}
\end{table}

\resultfig{task3_speedup}{Task 3: speedup of SPP, Dynamic I-RRIP and their combination.}{fig:t3-speedup}
\resultfig{task3_accuracy_coverage}{Task 3: prefetch accuracy and coverage.}{fig:t3-acc}
\resultfig{task3_l1d_mpki}{Task 3: L1D load MPKI.}{fig:t3-l1d}
\resultfig{task3_l2_mpki}{Task 3: L2C, L2 instruction and L2 data load MPKI.}{fig:t3-l2}

\paragraph{Discussion.} \todo{Show that the combination outperforms both standalone SPP and Dynamic I-RRIP in IPC, and explain where the extra gain comes from.}

% ==========================================
% SUMMARY
% ==========================================
\section{Summary}\label{sec:summary}

\begin{table}[htbp]
\centering
\caption{Geometric-mean speedup of every configuration over the baseline.}\label{tab:summary}
\tablesetup
\begin{tabular}{lrrrr}
\toprule
\rowcolor{tablehead}\colhead{Configuration} & \colhead{Trace 1} & \colhead{Trace 2} & \colhead{Trace 3} & \colhead{Geomean} \\
\midrule
IP-stride               & \tbd & \tbd & \tbd & \tbd \\
SPP                     & \tbd & \tbd & \tbd & \tbd \\
SRRIP                   & \tbd & \tbd & \tbd & \tbd \\
I-RRIP                  & \tbd & \tbd & \tbd & \tbd \\
Dyn.\ I-RRIP            & \tbd & \tbd & \tbd & \tbd \\
SPP + Dyn.\ I-RRIP      & \tbd & \tbd & \tbd & \tbd \\
\bottomrule
\end{tabular}
\end{table}

\resultfig[0.85\linewidth]{summary_speedup}{Speedup of all configurations.}{fig:summary}

\todo{Short closing paragraph: key takeaways from each task.}

% ==========================================
% APPENDIX
% ==========================================
\appendix
\section{Reproducing the Results}\label{app:repro}
\begin{cppcode}[Build and run]
./build_champsim.sh no lru baseline
./build_champsim.sh spp dyn_irrip spp_dynirrip
./bin/<binary> -warmup_instructions 25000000 \
    -simulation_instructions 25000000 -traces Trace_Files/Trace1.gz
\end{cppcode}
\todo{Describe the run and result-parsing scripts in the repository and where the raw outputs are stored.}


\end{document}
